\section{The Verification Asymmetry}
\label{sec:verification_asymmetry}

This section isolates the mechanism that makes the max-aggregation
pathology structurally irresolvable within the existing risk
framework, by formalizing the two observation regimes induced by a
risk-restriction pair and showing that one of them produces
trivial Bayesian updates.

\subsection{Formalizing the asymmetry}

\begin{definition}[Risk-Restriction Pair]
\label{def:risk_restriction}
A \emph{risk-restriction pair} $(R, \Xirestrict)$ consists of:
\begin{itemize}
\item A \emph{risk claim}
  $R = (a_j, (d_1, \ldots, d_k), \Pathway, \pR)$, where $a_j$ is
  the claiming agent, $(d_1, \ldots, d_k)$ is the capability tuple
  under concern, $\Pathway$ is the predicted exercise
  pathway~\citep[Definition~2]{lasser2026horizon} that would
  produce $\volL$-contraction, and $\pR \in (0, 1]$ is $a_j$'s
  claimed probability that $\Pathway$ activates if the tuple is
  exercised.
\item A \emph{restriction} $\Xirestrict$: the preemptive
  restriction of a keystone capability $d_{j^*}$ in the tuple
  (in the sense
  of~\citep[Proposition~1]{lasser2026horizon}), which prevents
  exercise of $\Pathway$.
\end{itemize}
\end{definition}

\begin{definition}[Observation Regimes]
\label{def:obs_regime}
Under a risk-restriction pair $(R, \Xirestrict)$, the
\emph{observation regime} at time $t$ is one of:
\begin{itemize}
\item \emph{Restricted regime} ($\Xirestrict$ active): the
  capability tuple is restricted; $\Pathway$ cannot activate.
  The only observations available are indirect---behavioural
  signals from $a_j$, correlational evidence from related
  capabilities, and time-passage without incident (uninformative
  because the restriction is what prevents the test).
\item \emph{Exercise regime} ($\Xirestrict$ lifted): the
  capability tuple is exercised; $\Pathway$ may or may not
  activate.  Direct observation of activation or non-activation
  is available.
\end{itemize}
\end{definition}

\begin{proposition}[Asymmetric Evidence Under Restriction]
\label{prop:asym_evidence}
Let $\pR(t)$ denote the posterior probability that risk claim $R$
is valid at time $t$.  Under the restricted regime:
\begin{itemize}
\item[(a)] If the pathway were to activate (counterfactual
  contraction observed under a hypothetical exercise), the
  standard Bayesian update yields
  $$\pR(t+1) = \pR(t) \cdot
     \frac{\Pr(\text{contraction} \mid R \text{ valid})}
          {\Pr(\text{contraction})}$$
  with
  $\Pr(\text{contraction} \mid R \text{ valid})
     > \Pr(\text{contraction} \mid R \text{ invalid})$,
  which increases $\pR$.
\item[(b)] Under the restricted regime, no exercise occurs, so
  $$\pR(t+1) = \pR(t).$$
  The posterior is unchanged because the restriction prevents
  the observation that would update it.  The absence of
  contraction under restriction is consistent with both $R$
  being valid (the restriction correctly prevented it) and $R$
  being invalid (there was nothing to prevent).
\end{itemize}
Part~(b) is the \emph{Wamura asymmetry}: successful prevention is
observationally equivalent to unnecessary prevention.
\end{proposition}

\begin{proof}[Proof sketch]
Direct from the Bayesian update rule and the counterfactual
structure.  Under restriction, the likelihood ratio
$\Pr(\text{no contraction} \mid R \text{ valid}, \Xirestrict)
 / \Pr(\text{no contraction} \mid R \text{ invalid}, \Xirestrict)
 = 1$
because both hypotheses predict no contraction when the pathway
is blocked.  Likelihood ratio~$1$ means the posterior equals the
prior: no information is gained.  The (a) direction is the
standard Bayes-positive update; the asymmetry in (a) versus (b)
is what produces the Wamura ratchet.
\end{proof}

\begin{remark}[Connection to substrate-level asymmetry]
\label{rem:substrate_asym}
The verification-asymmetry definition
of~\citep[Definition~2]{lasser2026exo} is the substrate-level
version of this result: a digital agent can demonstrate
capability possession (positive claim) far more easily than an
external observer can verify capability absence (negative claim).
Proposition~\ref{prop:asym_evidence} is the risk-framework
specialization: a risk claim can be confirmed by observing
catastrophe but cannot be disconfirmed by observing
non-catastrophe-under-restriction.  The monotonicity
result~\citep[Proposition~1]{lasser2026exo} ensures the
substrate-level asymmetry is non-decreasing in capability scale;
Proposition~\ref{prop:asym_evidence} inherits the same monotone
character for the risk-framework projection.
\end{remark}

\subsection{The overcaution accumulation dynamic}

\begin{proposition}[Monotone Overcaution Under Max-Aggregation]
\label{prop:monotone_overcaution}
Under max-aggregation of risk estimates~\citep{lasser2026horizon}
and neutral-prior initialization:
\begin{itemize}
\item[(a)] For any newly discovered capability tuple $(d_1,
  \ldots, d_k)$ with no exercise history, there exists a
  threshold $p_{\mathrm{thresh}} > 0$ such that if any agent $a_j$
  assigns $\pR(0) > p_{\mathrm{thresh}}$, the tuple is
  immediately restricted, by the preemptive restriction
  criterion of~\citep[Proposition~1]{lasser2026horizon}.
\item[(b)] Once restricted, $\pR(t) = \pR(0)$ for all $t > 0$
  under the restricted regime, by
  Proposition~\ref{prop:asym_evidence}(b).
\item[(c)] The number of restricted capability tuples is
  monotonically non-decreasing over time: new tuples may be
  discovered and restricted, but no restricted tuple is
  unrestricted through the standard evidence-accumulation
  mechanism.
\end{itemize}
\end{proposition}

\begin{proof}[Proof sketch]
Part~(a) follows from the restriction criterion
of~\citep[Proposition~1]{lasser2026horizon}: it compares
$\Risk_j$ against the restriction cost, and for any tuple a
sufficiently high $\pR$ exceeds the threshold.  Part~(b) is
Proposition~\ref{prop:asym_evidence}(b).  Part~(c) is immediate
from (b): no evidence accumulates to reduce $\pR$ below
$p_{\mathrm{thresh}}$, and the existing framework contains no
mechanism that generates exercise-regime observations for
restricted tuples.
\end{proof}

\begin{corollary}[Asymptotic Lockup]
\label{cor:lockup}
In the limit of many capability discoveries, the fraction of the
capability poset under active restriction grows monotonically,
bounded below by the fraction of tuples for which at least one
agent assigned $\pR(0) > p_{\mathrm{thresh}}$.  If agents are
systematically overcautious ($\pR(0)$ exceeds the true risk on
many tuples), the locked fraction can grow without bound
relative to the truly risky fraction, producing $\volL$
contraction that the framework cannot distinguish from genuine
safety.
\end{corollary}

The corollary is the formal expression of the Fudai floodgate
intuition taken to its pathological limit: if every agent assigns
moderately high initial risk to every newly discovered
capability, and no mechanism in the framework is able to
differentiate genuine risk from overcaution, the poset contracts
under restriction without any mechanism for reversal.  This is
precisely the degenerate phase boundary of Loop~3
in~\citep[Proposition~3]{lasser2026phase}: the self-correction
rate $\rS$ is zero because no exogenous signal on the restricted
dimension exists.  The remainder of the paper constructs the
missing mechanism.
